\begin{textsample}{POS Dim 3 – ro – Score 34.00 – ro/ro\_inf\_39.txt}  \label{ex:f3_pos_131}
Para propor intencionalidades educativas significativas, docentes da primeira \textbf{infância} podem : Pensando nos \textbf{Bebês} - Construir espaços cuidadosamente planejados, que permitam exploração livre e ampliação da \textbf{percepção} espacial ao deslocar -se enfrentando obstáculos presentes nos trajetos : subindo, descendo, pulando, passando por cima, por baixo, rodeando, equilibrando -se, ao explorar diferentes caminhos para se chegar a um mesmo lugar, ao procurar objetos ou pessoas que estão \textbf{escondidas} em diferentes lugares. - Oportunizar a exploração de objetos com formas e \textbf{volumes} variados, algumas propriedades simples dos materiais, por exemplo, a luminosidade, a temperatura, a consistência, a \textbf{textura} e inclinação dos diferentes tipos de solo da unidade de Educação \textbf{Infantil}. - Possibilitar a exploração de alimentos, objetos e cheiros que ampliem suas experiências visuais, táteis, auditivas, gustativas e olfativas, bem como a produção e \textbf{manuseio} de massinhas caseiras, experimentando as sensações. - Organizar momento para que \textbf{brinquem} com materiais com possibilidades transformadoras : com água e \textbf{areia}, ou com terra, “ melecas ”, pasta de maisena ou outros materiais, bem como objetos para serem amassados ou deslocados. - Incentivar que apreciem e acompanhem corporalmente o canto da \textbf{professora} alterando o ritmo e o timbre ( alto, baixo, grave, agudo ) dos \textbf{sons}. - Propiciar que explorem os espaços externos da \textbf{instituição}, conhecendo o trajeto e o percurso.
Pensando nas \textbf{Crianças} bem \textbf{pequenas} e \textbf{pequenas} - Possibilitar a exploração de objetos de diferentes formatos e tamanhos e que utilizem o conhecimento de suas propriedades para explorá -los com maior intencionalidade : por exemplo, empilhar objetos do menor para o maior e vice-versa. - Mediar que realizem ações ( parar uma bola, fazer bolinhos de \textbf{areia}, arrumar formas de carregar objetos pesados etc. ) e expliquem o que usou, de que forma fez. - Possibilitar que resolvam problemas cotidianos, como a divisão de materiais coletivos, a escolha da bola mais leve, a execução de uma receita que envolve medidas etc., desenvolvendo noções relativas à direção, sentido, quantidade, tempo. - Oportunizar que modelem massinha produzida a partir de um mingau grosso de água e maisena, pesquisem algumas de suas características, por exemplo : consistência ( duro, mole ), temperatura ( quente, frio ), \textbf{peso} ( leve, pesado ). - Organizar e mediar observação de fenômenos e elementos da natureza presentes no \textbf{dia} a \textbf{dia} e o reconhecimento de algumas características do clima : calor, chuva, claro-escuro, quente-frio. - Possibilitar que explorem traços e formas utilizando os materiais e procedimentos do fazer \textbf{plástico}. - Oportunizar que observem as características físicas dos animais ( os \textbf{sons} por eles produzidos, sua pelagem, forma do corpo, presença de bico, localização dos olhos e outras ), além de alimentação e moradia. - Possibilitar a participação em atividades que envolvam processos de \textbf{culinária}, levantando questões relativas à transformação dos ingredientes usados. - Garantir que explorem quantidades nas \textbf{brincadeiras} e práticas cotidianas, \textbf{brinquem} de recitar os números nas \textbf{brincadeiras} tradicionais. - Oferecer a exploração das relações de \textbf{peso}, tamanho e \textbf{volume} de formas bidimensionais ou tridimensionais. - Possibilitar a exploração de materiais como a argila e massa de modelar, percebendo a transformação do espaço tridimensional em bidimensional, e vice-versa, a partir da construção e desconstrução. - Garantir a utilização de diferentes instrumentos de \textbf{medição} convencional e não convencional a fim de estabelecer : distâncias, comprimento, capacidade ( litro ) e massa. - Oportunizar que expliquem o efeito e a transformação na forma, velocidade, \textbf{peso} e \textbf{volume} de objetos, agindo sobre eles. - Oportunizar a exploração e utilização de instrumentos para as medidas de comprimento, de tempo e temperatura, como : réguas, trenas, fitas métricas, prumo, esquadro, nivelador, termômetro, utensílios elétricos, palmos, passos, barbantes, relógios digitais e analógicos, \textbf{calendários} anuais etc. - Produção de desenhos utilizando instrumentos tecnológicos e software como o Paint. - Possibilitar exploração de algumas propriedades dos objetos, como a de refletir, ampliar ou inverter as imagens, ou de produzir, transmitir ou ampliar \textbf{sons} etc. - Oportunizar que observem e criem explicações para fenômenos e elementos da natureza presentes no seu \textbf{dia} a \textbf{dia} ( calor produzido pelo sol, chuva, claro-escuro, quente-frio ), estabelecendo regularidades, relacionando -os à necessidade dos humanos por abrigo e \textbf{cuidados} básicos — agasalhar -se, não ficar exposto ao sol, beber líquido, fechar ou abrir janela, acender ou apagar a luz —, refletindo sobre algumas mudanças de \textbf{hábitos} em animais ou plantas, influenciadas por mudanças climáticas. - Garantir a exploração e comparação de altura dos \textbf{colegas}, \textbf{medição} de ingredientes em receitas \textbf{culinárias}, a distância de um salto. - Garantir exploração de diferentes contextos sociais em que a utilização de números e contagem sejam necessárias, utilizando diferentes estratégias. - Garantir comunicação de quantidades a partir da linguagem oral e de \textbf{registros} escritos espontâneos de números em situações contextualizadas. - Possibilitar que solucionem problemas cotidianos na \textbf{instituição} relativos a noções geométricas, numéricas, espaciais e de medidas : cálculo de idades, altura, número de gols, datas, pontos de um jogo, entre outros. - Garantir a participação em jogos de regras com tabelas para \textbf{registro} e \textbf{acompanhamento} dos pontos obtidos. - Favorecer \textbf{brincadeiras} que possam utilizar mapas para solucionar problemas, como procurar objetos ou pessoas em diferentes lugares, buscando a posição deles. - Possibilitar que desenhem ou interpretem imagens de objetos a partir de diferentes pontos de vista ( de frente, de cima, de lado ). - Oportunizar e mediar a investigação das transformações de misturas, como a de água e \textbf{areia}, e outros elementos cotidianos, observando e experimentando as diferenças. - Possibilitar que planejem a organização do espaço da sala referência, utilizando para tanto o desenho para projetar tal organização. - Oportunizar que observem e comentem obras de artes visuais que exploram formas simétricas. - Oportunizar a utilização de materiais para construir imagens e objetos em espaços bidimensionais e tridimensionais. - Garantir o jogo simbólico com materiais que \textbf{convidem} a pensar sobre os números, como \textbf{brincar} de comprar e vender, identificando notas e moedas do sistema monetário vigente. - Oportunizar observação e estabelecimento de relações de diferença e de igualdade entre espécies vegetais. - Mediar e possibilitar que pesquisem a localização — em uma régua, fita métrica, quadro numérico ou \textbf{calendário} — buscando solucionar problemas cotidianos. - Possibilitar que pesquisem sobre os números em diferentes contextos sociais, como a numeração das casas da rua ; a localização do número dentro de uma coleção do interesse das \textbf{crianças} etc. - Garantir que explorem as notações numéricas em diferentes contextos : \textbf{registro} de jogos, controle de materiais da sala, quantidade de \textbf{crianças} que vão merendar ou que vão a um \textbf{passeio}. - Possibilitar que percebam alterações ocorrendo em seu próprio corpo : a perda e aparecimento de dentes, aumento na altura, no tamanho das \textbf{mãos} e dos \textbf{pés}, entre outras. - Oferecer e organizar momento para que identifiquem algumas características do ambiente e/ou das pessoas em \textbf{fotos}, relatos e outros \textbf{registros}, apontando semelhanças e diferenças com o tempo presente. - Propiciar que identifiquem a passagem do tempo \textbf{apoiadas} no \textbf{calendário} e utilizando a unidade de tempo — \textbf{dia}, \textbf{mês} e ano — para marcar as datas significativas para o grupo.

% matched lemmas: acompanhamento, apoiar, areia, bebê, brincadeira, brincar, calendário, colega, convidar, criança, cuidado, culinário, dia, esconder, foto, hábito, infantil, infância, instituição, manuseio, medição, mão, mês, passeio, pequeno, percepção, peso, plástico, professora, pé, registro, som, textura, volume
\end{textsample}
